\subsubsection{Analisis STRIDE pada Sistem DecMed}
\label{subsec:analisis-stride-decmed}
Setelah kelemahan fungsional pada arsitektur DecMed diidentifikasi pada Subbab \ref{subsec:analisis-sistem-decmed}, analisis STRIDE pada subbab ini digunakan sebagai pelengkap untuk memetakan implikasi keamanannya secara lebih sistematis. Dengan kata lain, STRIDE tidak digunakan untuk merumuskan masalah penelitian yang baru, melainkan untuk menunjukkan bahwa keterbatasan pada kontrol akses dan mekanisme delegasi memang berpotensi menimbulkan ancaman keamanan yang nyata ketika sistem diterapkan. Berdasarkan arsitektur DecMed pada Gambar \ref{fig:decmed-archi}, \textit{attack surface} utama sistem mencakup \textit{Patient Client}, \textit{Hospital Client}, \textit{Ministry Client}, data pada lapisan \textit{on-chain}, data terenkripsi pada lapisan \textit{off-chain}, \textit{instance} PRE, serta token kapabilitas yang mengatur otorisasi akses.

\textbf{\textit{Spoofing}.} Ancaman \textit{spoofing} muncul ketika penyerang berhasil menyamar sebagai pasien, dokter, atau institusi yang sah untuk meminta akses terhadap RME. Pada DecMed, risiko ini relevan karena keputusan akses sangat bergantung pada identitas pihak yang memegang token kapabilitas. Jika identitas klien atau kepemilikan token tidak dapat diverifikasi dengan baik, maka kredensial yang dicuri atau sesi yang dibajak dapat digunakan untuk memperoleh akses tidak sah terhadap data pasien.

\textbf{\textit{Tampering}.} Ancaman \textit{tampering} dapat menargetkan metadata rekam medis pada lapisan \textit{on-chain}, data terenkripsi pada IPFS, maupun token akses yang beredar antarpihak. Walaupun penyimpanan \textit{on-chain} membantu menjaga integritas referensi data, sistem tetap menghadapi risiko modifikasi tidak sah pada metadata akses, penggantian objek \textit{off-chain}, atau penyalahgunaan token yang tidak membawa batasan cukup rinci. Ancaman ini memperjelas bahwa keterbatasan representasi hak akses dan organisasi data pada DecMed tidak hanya menjadi persoalan desain, tetapi juga berdampak langsung pada integritas dan ketepatan pemberian akses.

\textbf{\textit{Repudiation}.} Ancaman \textit{repudiation} berkaitan dengan kemungkinan pihak tertentu menyangkal pernah melakukan permintaan, pemberian, atau penggunaan akses. Pada sistem RME, jejak audit penting untuk membuktikan siapa yang mengakses data, kapan akses dilakukan, dan data apa yang diproses. Walaupun aspek ini bukan fokus utama solusi yang dirancang dalam penelitian ini, pemetaan STRIDE menunjukkan bahwa mekanisme akses yang melibatkan penerbitan, delegasi, dan penggunaan token tetap membutuhkan penelusuran yang memadai agar penyalahgunaan akses dapat diinvestigasi.

\textbf{\textit{Information Disclosure}.} Ancaman ini merupakan salah satu risiko paling kritis pada DecMed karena sistem mengelola data kesehatan yang bersifat sensitif. Kebocoran dapat terjadi tidak hanya pada isi data klinis \textit{off-chain}, tetapi juga pada metadata \textit{on-chain} yang dapat mengungkap relasi pasien, fasyankes, atau jenis layanan yang diterima. Dalam konteks penelitian ini, ancaman \textit{information disclosure} secara langsung menguatkan kebutuhan akan kontrol akses yang lebih granular dan segmentasi data \textit{off-chain}, karena tanpa keduanya pemberian akses berpotensi membuka cakupan data yang lebih luas daripada kebutuhan aktual tenaga medis.

\textbf{\textit{Denial of Service}.} Ancaman \textit{denial of service} dapat diarahkan ke bagian dari DecMed yang dibutuhkan seperti layanan yang memproses akses PRE, jaringan DLT IOTA, dan IPFS. Meskipun pendekatan terdesentralisasi membantu mengurangi risiko \textit{single point of failure}, ketergantungan pada beberapa komponen terpisah tetap membuka kemungkinan layanan menjadi lambat atau tidak tersedia. Ancaman ini menunjukkan bahwa ketahanan operasional tetap penting pada sistem RME terdesentralisasi, meskipun tidak menjadi fokus utama perancangan mekanisme kontrol akses pada penelitian ini.

\textbf{\textit{Elevation of Privilege}.} Ancaman \textit{elevation of privilege} terjadi ketika pengguna yang semula hanya memiliki hak terbatas berhasil memperoleh hak akses yang lebih tinggi. Pada DecMed, risiko ini berkaitan langsung dengan model token kapabilitas yang masih statis dan belum mendukung atenuasi kewenangan secara rinci. Oleh karena itu, ancaman ini memperkuat kebutuhan akan mekanisme delegasi yang dapat membatasi jenis hak, ruang lingkup data, dan konteks penggunaan secara bertahap agar eskalasi hak akses di luar kebutuhan klinis yang sah dapat dicegah.

Secara keseluruhan, hasil analisis STRIDE ini menegaskan bahwa kelemahan yang telah diidentifikasi sebelumnya pada DecMed memang memiliki konsekuensi keamanan yang nyata. Di antara berbagai kategori ancaman, \textit{information disclosure}, \textit{tampering}, dan \textit{elevation of privilege} merupakan ancaman yang paling langsung berkaitan dengan fokus penelitian ini, yaitu kebutuhan akan kontrol akses granular, segmentasi data \textit{off-chain}, dan delegasi berantai yang terbatas. Dengan demikian, STRIDE berfungsi sebagai analisis pendukung yang memperkuat urgensi permasalahan, sedangkan perumusan masalah penelitian tetap diturunkan terutama dari keterbatasan fungsional DecMed, kebutuhan layanan kesehatan kolaboratif, dan tuntutan regulasi.
